% Source: https://artofproblemsolving.com/community/c3h3535499_construct_problem
% Stable identity: aops-topic-3535499-q01
Construct $\sqrt5$ as the hypotenuse of a right triangle with legs $1,2$, then construct $\phi=\dfrac{1+\sqrt5}{2}$ by addition and bisection. A right triangle with hypotenuse $1$ and adjacent leg $\dfrac{\phi}{2}$ constructs an angle $\theta$ with $\cos\theta=\dfrac{\phi}{2}$. To verify $\theta=36^\circ$, use $\phi^2=\phi+1$ to obtain $\cos3\theta=\dfrac{1-\phi}{2}=-\cos2\theta=\cos(180^\circ-2\theta)$. Since $0<\theta<45^\circ$, both angles lie between $0^\circ$ and $180^\circ$, where cosine is one-to-one. Hence $3\theta=180^\circ-2\theta$, giving $\theta=36^\circ$. Bisect it to construct $18^\circ$. Construct $60^\circ$ using an equilateral triangle and bisect twice to get $15^\circ$. Copy the $15^\circ$ and $18^\circ$ angles onto the same initial ray and same side. The angle between their other rays is $\boxed{3^\circ}$. Every operation described is a compass-and-straightedge construction.
\begin{asy}
size(7cm);
usepackage("amsmath");
real phi=(1+sqrt(5))/2;pair O=(0,0),A=(phi/2,0),B=(phi/2,sqrt(1-phi*phi/4));draw(O--A--B--cycle);label("$O$",O,SW);label("$A$",A,SE);label("$B$",B,NE);label("$\dfrac{\phi}{2}$",(O+A)/2,S);label("$1$",(O+B)/2,NW);label("$36^\circ$",(.18,.05),NE);
\end{asy}

\begin{asy}
usepackage("amsmath");
size(7cm);
pair O=(0,0);draw(O--5*dir(0));draw(O--5*dir(15));draw(O--5*dir(18));label("$O$",O,SW);label("$15^\circ$",5*dir(15),SE);label("$18^\circ$",5*dir(18),NE);draw(arc(O,3.5,15,18),blue);label("$3^\circ$",3.5*dir(16.5),N);
\end{asy}
